\section{The Categorical Data Model}
\label{sec:data-model}

\subsection{Basic Definitions: The Astrolabe Category}
\label{sec:dm-definitions}

% Definition 3.1.1 (Astrolabe Knowledge Graph).
% G = (O, M, s, t, σ_O, σ_M, h) where:
% O: objects, M: morphisms, s,t: source/target maps
% σ_O, σ_M: sort classification functions, h: unique identifier (Hex₁₂)
% Each o ∈ O carries optional rich-text: name, statement, proof, intuition, notes (MDX)
% Each m ∈ M carries optional: sort, notes

% Definition 3.1.2 (Astrolabe Category).
% A(G) = free category generated by G
% Objects are O, morphisms are finite composable paths from M (+ identity morphisms)
% Composition exists and is associative, sort propagates along composition
% Explicit morphisms are generators; composed morphisms are implicit (computed dynamically)
% Design is deliberate: storing all compositions → combinatorial explosion

% Definition 3.1.3 (Equivalence Relations and Commutative Diagrams).
% Users can declare path equivalences ~, forming quotient category A(G)/~
% Commutative diagrams = path equivalences (future work, A(G) is pure free category for now)

% Definition 3.1.4 (Plugin Functors).
% Import plugin = functor F: E → A(G) from external data source category
% F_Lean: L → A(G) maps Lean environment category to Astrolabe category
% Analysis plugin = endofunctor Φ: A(G) → A(G), enrichment (attaches metadata)

% Two-layer identifier design:
% Core layer: h = random 12-hex id (stable graph-internal references)
% Formalization plugin layer: content-addressed hash h_c from type signature (enables cross-instance merging)

% Implementation correspondence: ilean parser = functor F_Lean
% Interoperability: import functor output is standard obj/mor, immediately relatable to any existing node

% Definition 3.1.5 (Inclusion Morphisms and Containment Structure).
% sort: "inclusion" → A belongs to B
% Many-to-one, one-to-many, transitivity, DAG constraint (no cycles)
% Rendering: Canvas nesting/clustering, ObjBlock collapsible children
% Difference from "uses": inclusion = part-of, uses = logical dependency

% Definition 3.1.6 (Core Morphism Vocabulary). 12 core mor sort types:
% I. Structural: (1) inclusion, (2) formalizes
% II. Logical: (3) uses, (4) proof-uses, (5) implies, (6) equivalent, (7) contradicts
% III. Semantic: (8) generalizes, (9) specializes, (10) annotates, (11) motivates
% IV. Meta: (12) cites
% [Main text covers 4 most important; full definitions in Appendix F]
% Vocabulary is open --- users can create arbitrary sort strings with autoColor

% Definition 3.1.7 (Core Object Vocabulary).
% Table: obj sort → semantics → status tracking → formalization relation
% definition, theorem, lemma, proposition, corollary, conjecture, example, remark, axiom
% lean-theorem, lean-definition, lean-lemma, lean-instance (auto-created by ilean plugin)
% obj sort also open --- users can create arbitrary sorts

\subsection{Why Category Theory}
\label{sec:dm-why-ct}

% Not hype, but precise engineering requirements:
% - Morphisms are first-class citizens (relations have content, not just lines)
% - Morphisms are classifiable (sort system distinguishes dependency types)
% - Hierarchical structure is a special case of typed mor (inclusion), no separate folder/tree mechanism needed
% - Morphisms are composable (transitive dependency inference)
% - Unified interface (NL nodes and formalized nodes use identical obj/mor structure)
%
% Core design advantage: sort system reduces "inventing new relationship types" from
% "change code add feature" to "fill in a string". Traditional tools need separate mechanisms
% for each relationship type (folders for hierarchy, [[link]] for references, \uses for dependencies).
% In Astrolabe, all are mor with different sort values.

\subsection{Comparison with Existing Data Models}
\label{sec:dm-comparison}

% Table: Astrolabe vs leanblueprint vs LeanArchitect vs Obsidian vs Stacks Project
% Dimensions: node typed sort, edge typed sort, edges as first-class,
% content-addressing, non-formalized content, formalization import, network analysis, language agnostic

\subsection{Sort System and Automatic Color Assignment}
\label{sec:dm-sort-system}

% Sort is free string, not enum --- key to domain agnosticism
% Automatic color: deterministic FNV-1a hash → HSL color wheel, built-in math sort presets
% obj sort and mor sort independent but share color assignment mechanism

\subsection{MDX Replacing LaTeX}
\label{sec:dm-mdx}

% LaTeX is typesetting language, not structural language
% LaTeX/TikZ fatal flaws in AI era: unreliable LLM generation, not general enough, PDF is closed terminal state
% MDX = Markdown + JSX, natively supports componentization
% MDX embeds LaTeX formulas (KaTeX/MathJax), math expressiveness preserved
% MDX embeds React components for interactive mathematical content:
%   - Geometric objects can rotate/zoom, variational processes can animate
%   - Convergence behavior interactive, proof steps progressively unfold
% Relation to executable papers (Koop, Silva, Freire 2011)
% AI compatibility: LLMs generate MDX more easily than structured LaTeX

\subsection{The Closed-Loop Pipeline}
\label{sec:dm-pipeline}

% [One figure + 1-2 paragraphs. Summarize the four phases as an application of the data model,
%  not a standalone methodology section.]
%
% Figure: NL → obj/mor network → network analysis → directed formalization → .ilean reflux → update → re-analyze
%         (closed loop arrow back to network analysis)
%
% Para 1: The categorical data model enables a closed-loop pipeline for mathematical formalization.
%   Phase 1: users structure natural language content as obj/mor (MDX + AI skills).
%   Phase 2: analysis plugins compute centrality/community/bottleneck signals on the typed network.
%   Phase 3: signals guide which nodes to formalize (human or AI).
%   Phase 4: formalization artifacts (.ilean) flow back via import functor F_Lean as standard obj/mor,
%   the network updates, and analysis re-runs. This is possible because the import functor produces
%   obj/mor indistinguishable from manually created ones --- they share sort, hash, and rendering.
%
% Para 2: The pipeline is an application demonstration, not a contribution in itself. The contribution
%   is the data model (Section 3.1-3.5) and the system (Section 4) that make it possible. Detailed
%   per-phase workflows are documented in the project's user guide, not in this paper.

\subsection{Three-Layer Coexistence Design}
\label{sec:dm-coexistence}

% [Academic content from former Section 4.5. This has real modeling substance.]
%
% Layer 1: Many-to-many formalization mapping (data model level).
%   One NL obj → multiple Lean obj (theorem split into statement + lemma + helper).
%   One Lean obj → multiple NL obj (shared lemma across theorems).
%   Typed mor (sort: "formalizes") naturally supports many-to-many, no extra mechanism.
%   Contrast: leanblueprint's \lean{decl_name} is one-to-one binding.
%
% Layer 2: Sort-driven differential rendering (visual level).
%   Different mor sorts trigger different visual representations on Canvas:
%   inclusion → nesting, uses → solid arrow, proof-uses → dashed, formalizes → double-line,
%   contradicts → red. Same obj/mor data model; sort value determines rendering mode.
%   Rendering rules are plugin-extensible.
%
% Layer 3: MDX components as knowledge reference DSL (content expression level).
%   Basic: <ObjBlock>, <ObjRef>, <MorBlock>, <MorRef> — live references into the knowledge graph.
%   Composite (designed): <NodeCluster>, <ProofTree>, <FormalizationStatus>, <DependencyDiff>.
%   AI can generate these components on demand.
%   Analogy: obj/mor = data types, sort = type system, MDX components = syntax,
%   plugins = standard library, AI = code generator.
